\section{Anbindung externer Datenquellen}
Zur Erweiterung der Funktionalität werden externe Datenquellen eingebunden.  
Die Kommunikation erfolgt über eine \gls{wlan}-Verbindung.  
Der ESP32 verbindet sich dabei mit dem lokalen Netzwerk und ermöglicht anschließend die Kommunikation mit externen Diensten.

\begin{lstlisting}[style=arduino,caption={Funktion zum Verbindungsaufbau mit dem \gls{wifi}-Access-Point.}]
void connectToWiFi() {
  WiFi.mode(WIFI_STA);
  WiFi.begin(WIFI_SSID, WIFI_PASSWORD);

  while (WiFi.status() != WL_CONNECTED) {
    bootModeUpdate();
  }
  Serial.println("WLAN verbunden.");
}
\end{lstlisting}

Die Zeit wird über \gls{ntp} synchronisiert.  
Zusätzlich werden Wetterdaten von der Web-\gls{api} Open-Meteo geladen.  
Die Datenübertragung erfolgt über \gls{https}, wobeidie Verarbeitung der Antworten im \gls{json}-Formaterfolgt.
Die zentrale Funktion zur Aktualisierung der Wetterdaten erzeugt zunächst die Anfrage-URL, ruft anschließend die Daten ab und liest daraus die Temperatur, den Wettercode und die Tageszeit aus.

\begin{lstlisting}[style=arduino,caption={Auszug aus der Wetterdatenabfrage.}]
String url = "https://api.open-meteo.com/v1/forecast?latitude=";
url += String(WEATHER_LAT, 4);
url += "&longitude=";
url += String(WEATHER_LON, 4);
url += "&current=temperature_2m,weather_code,is_day";
int httpCode = http.GET();
String payload = http.getString();
deserializeJson(doc, payload);
outsideTemperature = doc["current"]["temperature_2m"] | 0.0;
weatherCode = doc["current"]["weather_code"] | -1;
isDay = doc["current"]["is_day"] | 1;
\end{lstlisting}

Die empfangenen Wetterdaten werden anschließend klassifiziert.  
Hierzu werden Temperaturbereiche und Wettercodes den definierten Wettergruppen zugeordnet.
Die Anzeige erfolgt daraufhin über entsprechende Wort- und Statusanzeigen auf den LED-Streifen.
Weitere Funktionen zur Wetterklassifikation, periodischen Aktualisierung und Fehlerbehandlung sind im vollständigen Programmcode in Anhang~\ref{apx:1} dokumentiert.